% sections/background.tex -- Quotations checked against the RFC texts fetched from rfc-editor.org
% on 2026-11-23 (RFC 9449, 8705, 9440, 9421) and vendor docs fetched the same day (W2/W10).
% Alg. 1 is transcribed from ../harness/verifier/core.py (verify()), in its check order (W25).
\section{Background}\label{sec:background}

\paragraph{DPoP.} RFC~9449 binds an access token $T$ to a client key $k$
through a confirmation claim \jkt and requires a fresh signed proof JWT with
each request~\cite{rfc9449}. The proof carries \htm, ``the value of the HTTP
method'', and \htu, ``the HTTP target URI \ldots{} without query and fragment
parts'', plus \ath, a hash of $T$. The narrowness is deliberate: ``the idea is
to sign just enough of the HTTP data to provide reasonable proof of possession
\ldots{} to avoid the substantial difficulties inherent in attempting to
normalize HTTP messages'' (\S4.2), and ``DPoP does not ensure the integrity of
the payload or headers of requests'' (\S11.7). The receiving server must
ensure that ``the \htm claim matches the HTTP method of the current request''
and that ``the \htu claim matches the HTTP URI value for the HTTP request in
which the JWT was received, ignoring any query and fragment parts''; servers
``SHOULD employ syntax-based normalization \ldots{} and scheme-based
normalization'' before comparing (\S4.3). What ``the current request'' means
when the checking server is not the executing server is not specified.
Replay protection is time-bounded; a single-use \jti check ``may not always
be feasible in practice, e.g., when multiple servers behind a single endpoint
have no shared state'' (\S11.1). Algorithm~\ref{alg:verifier} gives the
check as our reference verifier implements it.

\begin{algorithm}[t]
\caption{RFC~9449 resource-server check of the reference verifier, on the
view $\view(r)=(m_v,u_v)$ an adapter reconstructs.}\label{alg:verifier}
\KwIn{view $(m_v,u_v)$; headers \hdr{Authorization}, \hdr{DPoP}; config $(\canon, W, \mathit{cache})$}
\KwOut{\KwAllow with $(\mathit{sub},\jkt)$, or \KwDeny with a reason}
require exactly one \hdr{Authorization: DPoP} $T$ (no Bearer) and one \hdr{DPoP} $\pi$\;
parse $\pi$; require \clm{typ}$=$\clm{dpop+jwt}, asymmetric \clm{alg}, public \clm{jwk}\;
\lIf{$\neg\mathsf{Verify}(\pi.\mathit{jwk}, \pi)$}{\KwDeny \texttt{bad\_signature}}
require claims \jti, \htm, \htu, \iat\;
\lIf{$\pi.\htm \neq m_v$}{\KwDeny \texttt{htm\_mismatch}\Cmt*[f]{\S4.3(8)}}
\lIf{$\canon(\pi.\htu) \neq \canon(u_v \setminus \{\mathit{query},\mathit{fragment}\})$}{\KwDeny \texttt{htu\_mismatch}\Cmt*[f]{\S4.3(9)}}
\lIf{$|\mathit{now}-\pi.\iat| > W$}{\KwDeny \texttt{iat\_skew}}
verify $T$ (signature, \clm{typ}, \clm{iss}, \clm{aud}, \clm{exp})\;
\lIf{$\pi.\ath \neq H(T)$}{\KwDeny \texttt{ath\_mismatch}\Cmt*[f]{\S7}}
\lIf{$T.\jkt \neq \mathsf{th}(\pi.\mathit{jwk})$}{\KwDeny \texttt{jkt\_mismatch}}
\lIf{$\mathit{cache}$ on $\wedge\ \pi.\jti \in \mathit{cache}$}{\KwDeny \texttt{jti\_replay}}
add $\pi.\jti$ to $\mathit{cache}$\Cmt*[f]{only on allow}
\KwAllow; emit \hdr{X-Verified-Sub}, \hdr{X-Verified-Jkt} (and \hdr{X-Verified-View}, Alg.~\ref{alg:db})\;
\end{algorithm}

\paragraph{Certificate binding and signatures.} RFC~8705 binds tokens to a
client certificate; behind a TLS-terminating intermediary, how certificate
metadata reaches the application ``is out of scope'' (\S6.5)~\cite{rfc8705}.
RFC~9440 standardises \hdr{Client-Cert} for this hop and requires that
incoming copies ``MUST be removed or overwritten before forwarding''~\cite{rfc9440}.
A certificate binding thus becomes a header whose soundness depends on every
hop. HTTP Message Signatures let the signer choose the covered components;
uncovered portions ``are susceptible to modification by an attacker without
affecting the signature'' (RFC~9421 \S7.2.1)~\cite{rfc9421}.

\paragraph{Layered verification.} Gateways rarely verify DPoP natively; they
delegate to an external service that reconstructs the client's request from
what the gateway sends. Envoy \texttt{ext\_authz} sends a check request with
the method and path~\cite{w2:envoyextauthz}; Traefik ForwardAuth sends
\hdr{X-Forwarded-Method}/\hdr{-Uri} and grants access on any
2xx~\cite{w2:traefikfwdauth}; NGINX \texttt{auth\_request} ``implements client
authorization based on the result of a subrequest'', and the original URI
reaches the service only through headers set in the configuration,
conventionally \hdr{X-Original-URI} from
\texttt{\$request\_uri}~\cite{w2:nginxauthreq}. Kong, APISIX and HAProxy have
equivalent mechanisms (\S\ref{sec:harness}). In each, the verifier sees one
rendering of the request while the router acts on another.
